\thispagestyle{empty}
\null
\vspace*{\fill}
\begin{flushleft}
\begin{minipage}{0.78\textwidth}
This project ships under a permissive license on purpose: teachers, students, and curious readers can treat the manuscript as a workshop that keeps breathing instead of a shrink-wrapped product. That is the same instinct that dragged academic software toward open repositories long before music pedagogy caught the habit. The prose, layout instructions, and engraved examples all move together as \gls{plain-text} sources tracked with \gls{version-control}, so a thoughtful correction lands with a paper trail instead of a mystery attachment.

You never owe the repo a pull request to benefit from it. Read quietly, steal a page for a lesson, or let a friend walk you through a passage. If you \emph{do} propose a change, the same openness lets someone fold it into the next revision with credit and care.
\end{minipage}
\end{flushleft}
\vspace*{\fill}
\clearpage
